%% ============================================================================
%%  subdocumento-ejemplo/contenido.tex
%%
%%  Ejemplo minimo de un subdocumento autocontenido. Sirve para verificar que
%%  el mecanismo funciona y como molde para los subdocumentos reales.
%%
%%  Reglas de la carpeta:
%%   - Este archivo SIEMPRE se llama contenido.tex y empieza por \chapter.
%%   - Las figuras van en figuras/ y se llaman por ruta relativa.
%%   - Los fragmentos van en tablas/ o donde quieras, y se traen con \parte{}.
%%   - NADA de aqui conoce la ruta de montaje: no escribas rutas que suban
%%     con ../ ni que empiecen desde la raiz del repositorio. Si lo haces,
%%     mover la carpeta rompe el capitulo.
%% ============================================================================

\setcounter{chapter}{7}
\chapter{Introducción a riesgos}


El plan establece tratamiento de riesgos de solución, desarrollo e implantación, con responsables, cuantificación y reservas durante el contrato.

La solución de LafroX integra recepción y trazabilidad, inventario, frío, preventa, reparto y cobranza en E1, y canal moderno, portales y costo de servir en E2. Sus riesgos dependen de una operación que despacha 96 camiones entre 05:30 y 07:00, mantiene los CD 24 horas sin WAN y el terreno 14 horas sin señal, conserva ERP como único emisor tributario y despliega sin detener rutas (Distribuidora Puelche S.A., 2026c, Cap. 10). Este plan conecta las decisiones comerciales del SD2, alcance del SD3, contratos y dimensionamiento del SD4, gobierno del SD6 y programación del SD7 (LafroX, 2026).

El cuerpo establece método y decisiones; los Anexos 8.A\textendash{}8.G complementan el análisis con fichas de riesgo, insumos de cuantificación y evidencia de planificación; el T-16 conserva el formato obligatorio separado (Distribuidora Puelche S.A., 2026d, sección 11). Las condiciones que deben cerrarse antes de cada hito, incluida la verificación del RPO y del riesgo residual que el SD4 declara en la sección 4.3.2.4, se registran en el Anexo 8.E con responsable y plazo.

\section{Plan de riesgos}

Esta sección define cómo se gestionan los riesgos durante los 56 meses: el ciclo y sus instancias, los roles y su autoridad, las escalas con que se analizan y las reglas de registro y cierre.

\subsection{Enfoque y ciclo}

El registro vivo sigue identificación, análisis, respuesta y seguimiento conforme a la norma ISO 31000 (International Organization for Standardization [ISO], 2018), que exige el RT-19.04 de las Bases Técnicas Transversales (Distribuidora Puelche S.A., 2026b). Cada entrada distingue causa, evento incierto y consecuencia, sin convertir una inconsistencia observada en riesgo futuro. El Anexo 8.E separa problemas y dependencias actuales.

JP mantiene el registro único. Los líderes identifican riesgos al revisar interfaces, pruebas, capacidad y despliegues; verifican disparadores semanalmente y ante cada cambio. El Comité de Proyecto quincenal revisa todos los riesgos abiertos, responsables, plazos y evidencia. Los comités Ejecutivo y de Arquitectura mensuales deciden escalamiento y cambios de sus ámbitos, y desde el mes 13 el Comité de Operación mensual revisa los riesgos de servicio: niveles de atención, incidentes, capacidad y continuidad (SD6, sección 6.1.6). Una amenaza a continuidad, seguridad o hito se escala inmediatamente. Marcha blanca exige seguimiento diario; operación mantiene vigilancia diaria de incidentes y colas y revisión en comités. En cada comité se compara la reserva restante con el riesgo remanente (análisis de reserva), y al cierre de cada etapa, en el H5 y el H10, una auditoría de riesgos comprueba que el proceso funciona (PMI, 2017, p. 456). El Comité Ejecutivo recibe cada mes un extracto con los diez riesgos de mayor exposición; el registro completo queda en el repositorio del proyecto.

\subsection{Roles y autoridad}

La Tabla 8.1 distingue al responsable del riesgo, la familia de esfuerzo y su vigencia. Los responsables dirigen equipos; no ejecutan solos las HH. Las familias del T-15 se conservan al relevar liderazgo.

\begingroup
\renewcommand{\thetable}{8.1}
\begin{tablalafrox}{Responsabilidad y vigencia por familia. Fuente: SD1 Tabla 1.3, SD13 y T-16.}{tab:sd8-cuerpo-1}{>{\hsize=0.48404\hsize\linewidth=\hsize\RaggedRight\arraybackslash}X >{\hsize=1.00760\hsize\linewidth=\hsize\RaggedRight\arraybackslash}X >{\hsize=1.28064\hsize\linewidth=\hsize\RaggedRight\arraybackslash}X >{\hsize=1.22771\hsize\linewidth=\hsize\RaggedRight\arraybackslash}X}{\cab{Familia} & \cab{Hasta mes 21} & \cab{Meses 22\textendash{}56} & \cab{Autoridad/aprobación}}
JP & Alex Aravena & Alex Aravena, dirección de contrato & Registro/escalamiento \\
ARQ & Bastián Trejo & Trejo, Comité Arquitectura & Diseño/frontera DR \\
SEG & Álvaro Catalán & Álvaro Catalán & Seguridad/privacidad \\
DAT & Leandro Chamorro & Guillermo Castillo; equipo DAT & Miño/Calidad: parámetros \\
DES & Tomás Pérez & Guillermo Castillo; equipo DES & CAL: verifica cambios \\
CAL & Maximiliano Miño & Maximiliano Miño & Evidencia/eficacia \\
SRE & Guillermo Castillo & Guillermo Castillo & Capacidad/continuidad \\
IMP & Patricio Henríquez & Guillermo Castillo; equipo IMP & CAL: usuarios; SEG: privacidad \\
\end{tablalafrox}
\endgroup

Las responsabilidades de DAT e IMP se transfieren al frente F8 de Operación, como establece el SD13. Para que no dependan de una sola persona, desde el mes 22 Bastián Trejo, que sigue en el Comité de Arquitectura, es suplente de Guillermo Castillo en DAT y DES, y Maximiliano Miño lo es en IMP; ambos continúan hasta el mes 56 (SD1, Tabla 1.3). R8-27 requiere aprobación sanitaria de Calidad; R8-29, privacidad de Seguridad. La Contraparte Técnica acepta con acta; Operaciones autoriza cortes y Calidad del CLIENTE decide sobre los aspectos sanitarios. Su agenda se fija en mes 1 y se sigue con R8-12. Los cambios físicos siguen el gobierno de arquitectura.

\subsection{Escalas previas al análisis}

Para poder realizar un buen análisis de riegos, debemos de definir la probabilidad (P), el impacto (I) y la prioridad (P x I).

La probabilidad cualitativa (P) expresa el juicio ordinal sobre la posibilidad de ocurrencia del riesgo. Se asigna en una escala de \textbf{1 a 5} mediante las bandas de ponderación cualitativa $q$ de la tabla 8.2. Estas bandas no se aplican a $p_0$, la estimación inicial porcentual de ocurrencia que se muestra por separado en T-16 y se usa para el valor esperado en HH.

\begin{tablalafrox}{Escala cualitativa de probabilidad -- Fuente: elaboración propia}{tab:tabla-probabilidad}%
  {G{0.08} G{0.16} Y Y}%
  {\cab{Valor} & \cab{Banda cualitativa ($q$)} & \cab{Descripción} & \cab{Criterio cualitativo}}
  1 & $0\%\leq q<20\%$ &Muy baja & Poco probable que ocurra; condiciones excepcionales \\
  2 & $20\%\leq q<40\%$ & Baja & Puede ocurrir en casos aislados \\
  3 & $40\%\leq q<60\%$ &Media & Podría ocurrir algunas veces durante el proyecto \\
  4 & $60\%\leq q<80\%$ &Alta & Ocurrirá con frecuencia si no se mitiga \\
  5 & $80\%\leq q\leq100\%$ &Muy alta & Casi seguro de que ocurra \\
\end{tablalafrox}

Por otro lado, el impacto (I) representa el nivel de afectación que tendria el riesgo si llegara a materializarse. Se evalúa también en una escala de \textbf{1 a 5}, donde 1 representa un impacto insignificante y 5 un impacto muy alto.

\begin{tablalafrox}{Escala de impacto -- Fuente: elaboración propia}{tab:tabla-impacto}%
  {F{0.24} Y Y}%
  {\cab{Valor} & \cab{Descripción} & \cab{Criterio cualitativo}}
  1 & Insignificante & Impacto imperceptible en el proyecto \\
  2 & Menor & Afectación leve, fácil de controlar \\
  3 & Moderado & Impacto parcial en entregables o plazos \\
  4 & Alto & Afecta hitos principales o desempeño general \\
  5 & Muy alto & Compromete el éxito total del proyecto \\
\end{tablalafrox}

Por último, luego de definir valores de P e I, se obtiene el producto $P \times I$, que nos da el nivel de prioridad del riesgo. El resultado se interpreta según la tabla 8.4, que clasifica los riesgos según su importancia relativa

\begin{tablalafrox}{Matriz de prioridad -- Fuente: elaboración propia}{tab:tabla-prioridad}%
  {F{0.24} Y Y}%
  {\cab{PxI} & \cab{Nivel de prioridad}}
  1 - 4 & Baja \\
  5 - 9 & Media \\
  10 - 16 & Alta \\
  17 - 25 & Crítica \\
\end{tablalafrox}

La probabilidad ordinal y el impacto ordinal determinan la exposición cualitativa; el impacto en HH se mantiene separado.

\subsection{Registro y cierre}

El responsable actualiza fuente, estado, disparador, tratamiento, consumo de HH y evidencia. Estados: identificado, en tratamiento, control verificado o materializado. Materializar el evento abre un problema y conserva su ID. CAL verifica cierre técnico; JP registra la decisión; la aceptación contractual corresponde al CLIENTE. La puntuación residual efectiva exige eficacia comprobada; C.5 sólo compara hipótesis. Verificar un control no libera reserva automáticamente. Ninguna aceptación de riesgo exime requisitos obligatorios.

\section{Identificación y Análisis de Riesgos}

Esta sección identifica los riesgos con una RBS y los analiza cualitativamente mediante la exposición ordinal ($P\times I$) y cuantitativamente mediante el valor esperado de cada riesgo en horas-hombre. La simulación Monte Carlo del cronograma por actividad del SD7 queda pendiente de definir y validar sus insumos; en esta versión no se presentan resultados simulados.

\subsection{RBS y análisis separados}

La estructura de desglose de riesgos (EDR, o RBS por sus siglas en inglés) tiene como raíz «Riesgos del proyecto LafroX» y organiza los 45 riesgos del registro actualizado, identificados desde R-001 hasta R-045. La Figura~\ref{fig:sd8-rbs} los distribuye en siete ramas según la fuente o dependencia predominante: Solución y tecnología (14), Integración y legado (7), Operación del cliente (8), Personas y adopción (7), Contrato y licitación (4), Regulatorio y externo (2) y Empresa proveedora (3). Cada riesgo conserva un ID único y una ubicación en la EDR, aunque sus consecuencias alcancen varias áreas.

\parte{figuras/rbs}

La rama \textbf{Solución y tecnología} reúne los riesgos del funcionamiento y la continuidad de la solución: calidad de datos, reconexión, autonomía móvil, operación en frío, trazabilidad sanitaria, capacidad, recuperación, seguridad y dependencias tecnológicas. Incluye, por ejemplo, los conflictos y duplicados tras reconexión (R-005), el retiro sanitario sin trazabilidad completa (R-012) y la recuperación y el retorno que exceden RTO/RPO (R-015). También incorpora las innovaciones de reproducción de incidentes (INN-02, R-037) y estimación de vida útil (INN-03, R-038). Es la rama con más riesgos; esa cantidad describe la distribución del registro y no determina por sí sola su prioridad.

\textbf{Integración y legado} agrupa las dependencias de interfaces del ERP, telemetría de terceros, convivencia del WMS, certificación EDI, avisos, pagos y datos de ruteo; entre ellas están las interfaces del ERP no utilizables (R-001) y la certificación EDI dependiente de las cadenas (R-030). \textbf{Operación del cliente} reúne las condiciones de conectividad, infraestructura, disponibilidad de contrapartes, ventanas de intervención, instalaciones y control de los flujos físicos y de efectivo. Ejemplos son la caída simultánea de enlaces de Concepción (R-006), las contrapartes que no validan a tiempo (R-022) y las devoluciones y envases que alteran saldo o custodia (R-044).

\textbf{Personas y adopción} aborda la objeción sindical y la privacidad de geolocalización (R-017), la incorporación de conductores rotativos por los transportistas (R-018), la rotación y el aprendizaje en bodega nocturna (R-019), la capacitación y la pérdida del conocimiento de rutas (R-021). En esta rama también se ubican el seguimiento de stock posentrega de INN-01 (R-036) y la comprensión y privacidad de la hoja del almacenero de INN-05 (R-040). Así, la EDR distingue las dependencias humanas de las condiciones técnicas y operativas que pueden afectar el mismo proceso.

\textbf{Contrato y licitación} concentra la aceptación y marcha blanca (R-026), los cambios de alcance, los incentivos de costo de servir de INN-04 y la transferencia y salida. \textbf{Regulatorio y externo} cubre la evidencia verificable del POD y la guía (R-029), junto con los cambios de retención o normativa aplicable (R-042). \textbf{Empresa proveedora} reúne la suficiencia de dotación compartida entre etapas (R-025), la capacidad de mesa, NOC y soporte de sitio (R-034), y la estimación de costos variables y caja del proveedor (R-035).

Los temas de riesgo señalados en el Caso 02 (Distribuidora Puelche S.A., 2026c, Cap. 19) se pueden recorrer con estos códigos: interfaces del ERP (R-001), enlaces de Concepción (R-006), objeción sindical (R-017), conductores de terceros (R-018), rotación en bodega (R-019), conocimiento de rutas (R-021) y capacidad ante el despacho crítico o el peak (R-013). Su distribución entre ramas permite reconocer las distintas fuentes de riesgo sin reunirlas artificialmente bajo una sola fase del proyecto.

Los análisis de \textbf{solución}, \textbf{desarrollo} e \textbf{implantación} exigidos por el Formulario T-22 se mantienen como perspectivas de evaluación, no como ramas de esta EDR. Solución examina el comportamiento, la integridad y la continuidad durante el uso; Desarrollo, la construcción, integración y certificación de los productos; e Implantación, la habilitación, aceptación y adopción del servicio en el CLIENTE. Un mismo riesgo puede afectar más de una perspectiva: la certificación EDI (R-030), por ejemplo, condiciona tanto el desarrollo de la integración como su habilitación. El análisis transversal conserva el ID sin duplicar riesgos.

La figura presenta la clasificación del registro actualizado de 45 riesgos, utilizada también en el Formulario T-16. Los análisis y anexos que conservan la codificación anterior R8 corresponden a una versión previa del registro; sus valoraciones y resultados no se atribuyen automáticamente a los códigos R-001 a R-045. La EDR permite localizar las fuentes de riesgo, mientras que el registro detalla sus causas, consecuencias, responsables y respuestas.

\subsection{Análisis cualitativo de riesgos}

El análisis cualitativo del T-16 combina probabilidad e impacto ordinales en una escala de 1 a 5. Su producto ($P\times I$) expresa la exposición, que se interpreta con las bandas de prioridad definidas en SD8. La Tabla 8.4 presenta los diez riesgos con mayor exposición, ordenados de mayor a menor valor; en caso de empate, se prioriza el ID ascendente.

\begingroup
\renewcommand{\thetable}{8.4}
\begin{tablalafrox}{Valoración cualitativa vigente de riesgos. Fuente: T-16; bandas de prioridad SD8.}{tab:sd8-cuerpo-4}{>{\hsize=0.8\hsize\linewidth=\hsize\RaggedRight\arraybackslash}X >{\hsize=1.1\hsize\linewidth=\hsize\RaggedRight\arraybackslash}X >{\hsize=1.05\hsize\linewidth=\hsize\RaggedRight\arraybackslash}X >{\hsize=1.1\hsize\linewidth=\hsize\RaggedRight\arraybackslash}X >{\hsize=1.4\hsize\linewidth=\hsize\RaggedRight\arraybackslash}X}{\cab{ID} & \cab{Prob. ordinal} & \cab{Impacto ordinal} & \cab{Exposición ($P\times I$)} & \cab{Nivel de prioridad}}
R-012 & 3 & 4 & 12 & Alta \\
R-015 & 3 & 4 & 12 & Alta \\
R-005 & 2 & 4 & 8 & Media \\
R-006 & 2 & 4 & 8 & Media \\
R-011 & 2 & 4 & 8 & Media \\
R-013 & 2 & 4 & 8 & Media \\
R-014 & 2 & 4 & 8 & Media \\
R-016 & 2 & 4 & 8 & Media \\
R-001 & 3 & 2 & 6 & Media \\
R-003 & 3 & 2 & 6 & Media \\
\end{tablalafrox}
\endgroup

Entre los diez seleccionados, R-012 y R-015 tienen la mayor exposición (12, prioridad Alta). Les siguen seis riesgos con exposición 8 y prioridad Media, y dos con exposición 6 y prioridad Media. El corte en exposición 6 deja varios riesgos empatados; para mantener exactamente diez filas se incluyen R-001 y R-003, los primeros por ID ascendente. La prioridad es relativa a la escala cualitativa y no equivale a un impacto estimado en horas hombre.

\subsection{Análisis cuantitativo}

El valor esperado cuantifica el esfuerzo condicional de cada riesgo ponderado por su probabilidad inicial estimada $p_0$. Para el riesgo $i$, se calcula como $VE_i=p_{0,i}\times H_i^{\mathrm{cond}}$, donde $p_{0,i}$ se expresa entre 0 y 1 y $H_i^{\mathrm{cond}}$ es el impacto de esfuerzo si el evento ocurre, medido en horas-hombre (HH).

La suma de los valores esperados de los 45 riesgos es \textbf{2.058,480 HH}. El detalle individual, las probabilidades iniciales y los impactos condicionales se presentan en el Anexo 8.B. Las estimaciones son de planificación y deben validarse con los responsables de riesgo.

Este total es un valor esperado analítico; no representa una reserva, un percentil, un esfuerzo cierto ni un resultado de simulación. La exposición cualitativa continúa calculándose por separado con los niveles ordinales de probabilidad e impacto. En este punto no se informan resultados Monte Carlo.

\section{Plan de Acción a Riesgos}

Esta sección fija las respuestas a los riesgos analizados, las reservas que las financian en horas y las condiciones que deben cumplirse antes de declarar factible la propuesta.

\subsection{Respuestas y costo-beneficio técnico}

Cada ficha del Anexo 8.A define estrategia, responsable temporal, disparador, plazo, mitigación, contingencia, residual hipotético, riesgo secundario y evidencia. Se mitigan 29 amenazas, incluida R8-14; R8-12 y R8-17 se escalan, y R8-22 se acepta activamente con medición y recálculo de la dotación de la mesa. La subcontratación no traslada la obligación final de LafroX.

La Tabla C.5 compara las 32 respuestas en HH (PMI, 2017, pp. 442\textendash{}443). Los 35 paquetes de control únicos suman 6.768,00 HH ya programadas. Su costo se cuenta una vez, aunque un control proteja varios riesgos; los retornos individuales no se suman. Se descarta toda alternativa que incumpla sanidad, despacho, integridad de datos o una función obligatoria.

Si la probabilidad disminuyera un nivel, el registro conjunto residual sería 14.055,17 HH, frente a las 20.198,98 HH iniciales. Esta hipótesis no acredita eficacia ni permite liberar automáticamente la diferencia de 6.143,81 HH. CAL comprueba resultados y JP solicita cambios con evidencia, remanente y ventana de uso. La certificación paralela origina R8-31; el refuerzo de calidad, R8-32; la captura de trazas, R8-24; y la separación de equipos mantiene presión sobre R8-11.

\subsection{Reservas y cronograma}

La reserva de contingencia cubre los riesgos identificados: \textbf{20.198,98 HH}, el 9,88 \% de las 204.527 HH base de los paquetes. R8-11 y R8-18 concentran el 54,8 \% de ella. La capacidad protegida E1 sólo se descuenta cuando coinciden los meses 13\textendash{}20, la Etapa 1 y los perfiles DES/CAL; los paquetes de R8-02 y R8-04 terminan antes del mes 13 y los de R8-14 son trabajo de la Etapa 2, de modo que no se descuenta nada y las 3.072 HH quedan protegidas para correcciones futuras de la Etapa 1, sin prestarse a la Etapa 2.

La Tabla 8.5 resume la contingencia por período; la Tabla D.2 la detalla por mes y perfil. El T-15, sección 4.5, la refleja en el plan: compara, período a período, el peak de personas equivalentes adicionales con la dotación. Toda la contingencia cabe en la dotación declarada, salvo implantación en los meses 15 y 20, donde la extensión de la marcha blanca (R8-18) se cubre con la opción de hasta 8 personas adicionales del contrato de implantación (E8-14).

\begingroup
\renewcommand{\thetable}{8.5}
\begin{tablalafrox}{Contingencia por período, HH. Fuente: Anexo 8.D, Tabla D.2.}{tab:sd8-cuerpo-6}{>{\hsize=1.01370\hsize\linewidth=\hsize\RaggedRight\arraybackslash}X >{\hsize=0.98630\hsize\linewidth=\hsize\RaggedRight\arraybackslash}X}{\cab{Período} & \cab{HH de contingencia}}
Meses 1\textendash{}6 & 5.000,78 \\
Meses 7\textendash{}12 & 5.253,15 \\
Meses 13\textendash{}15 & 3.230,04 \\
Meses 16\textendash{}21 & 2.995,70 \\
Meses 22\textendash{}33 & 1.278,99 \\
Meses 34\textendash{}56 & 2.440,32 \\
Total & 20.198,98 \\
\end{tablalafrox}
\endgroup

La distribución mantiene reserva hasta el mes 56. Cada cargo identifica evento, paquete, etapa, perfil, mes, HH y evidencia para impedir el doble consumo. Sólo una misma corrección comparte cargo; el residual hipotético no se libera sin cierre y autorización.

La reserva de gestión es de \textbf{1.600 HH}, separada de la contingencia y fuera de la línea base. Cubre un evento que el registro no identifica, dimensionado como rehacer un módulo de clase D (960 HH de desarrollo), su integración (480 HH) y una certificación (160 HH), con las unidades del T-15, sección 4.1. Con 8 desarrolladores y 4 evaluadores adicionales toma 0,94 meses de desarrollo y 1,25 de integración y certificación, unos 2,2 meses en secuencia. La autoriza el Comité Ejecutivo y obliga a actualizar la línea base; un riesgo identificado nunca se carga a esta reserva.

La reserva de cronograma son los días hábiles entre la entrega programada de cada hito y su fecha límite, de 4 a 35 (T-15, Tabla 5.2). Las cuatro semanas finales de cada marcha blanca son evidencia obligatoria y no aportan reserva. La valorización y el flujo monetario de las reservas corresponden a la Oferta Económica (BA Art. 50.2).

\subsection{Factibilidad y aceptación}

El Anexo 8.E registra las condiciones de evidencia: productividad/dotación, secuencias diarias y plazos de revisión/subsanación del Art. 18.3, fecha contractual, continuidad de 96 despachos, RPO remoto, atención y evidencia láctea. La coordinación de reserva y retención de INT-03/04 (SD4, apartado 4.1.4.4 y Anexo 4-G) está dimensionada en el Anexo 4-I, Tabla A.10, y el Anexo 4-W, Tablas A.32 y A.33, con un máximo de cuatro mensajes por línea. La retención consumida no se libera y esa holgura cubre líneas repartidas entre lotes o sitios. AL-STOCK-01 mide la proporción de líneas con más de una retención antes del H4. Las pruebas 3.8.4 y 3.9.3 verifican la carga y el drenaje, sin sumar la coordinación al drenaje tras un corte. El compromiso RPO \ensuremath{\leq}15 min del SD4, sección 4.2.5.1, se verifica por escenario sobre fibra, LTE y Starlink. La falla simultánea de los tres caminos seguida de la destrucción del sitio antes de reponer alguno es el riesgo residual declarado en SD4 4.3.2.4 (RT-02.11). R8-05 lo trata con alarmas de retraso de replicación a los 5 y 15 min, reposición del enlace por el proveedor, preemisión de guías al cerrar la carga y conservación local en el NAS WORM de Talca. AL-DR-01 mide RPO/RTO antes de H5/H10. Si el NAS se destruye con el sitio, sólo queda la última copia remota consistente: las medidas locales no demuestran por sí solas una pérdida de datos de 15 minutos o menos. Por eso E8-05 exige demostrar en AL-DR-01 qué copia remota sobrevive y en qué frontera se cumple el RPO de 15 minutos (RT-07.04), junto con el tratamiento del riesgo residual (RT-02.11).

El retiro sanitario debe identificar a los clientes afectados con evidencia en menos de dos horas; los 85 minutos de diseño del SD5 no constituyen un ensayo aprobado. E8-13 y E8-14 establecen las condiciones de calendario y capacidad. No se autoriza corte sin continuidad medida ni aceptación sin las seis condiciones simultáneas del Art. 17.3 y acta según Art. 18. La suspensión del proveedor de lácteos, de marzo a septiembre de 2026, es antecedente ocurrido. La consulta V-13 del SD3 (Anexo 3.H) precisa en el mes 1, con el CLIENTE y el proveedor, qué evidencia de trazabilidad restablece la relación, y la Etapa 1 prioriza esa trazabilidad; la restitución depende del proveedor. La propuesta queda sujeta a las condiciones de cierre del Anexo 8.E, cada una con responsable, hito límite y evidencia; ninguna se da por cumplida sin las decisiones, pruebas y actas que allí se indican.

\section*{Referencias}
\addcontentsline{toc}{section}{Referencias}

Las fuentes citadas en este documento se listan en formato APA 7.ª edición. Las Bases se citan en el texto con su documento y el artículo, capítulo, sección o código del requisito.

\begin{itemize}[leftmargin=*,itemsep=3pt]
\item Distribuidora Puelche S.A. (2026a). \emph{Bases Administrativas de Licitación N.º TFEP-01/2026: Contratación de Solución Integral de Software y Servicios de Operación}.
\item Distribuidora Puelche S.A. (2026b). \emph{Bases Técnicas Transversales de Licitación N.º TFEP-01/2026}.
\item Distribuidora Puelche S.A. (2026c). \emph{Caso 02: Logística. Especificaciones del problema y operación de Distribuidora Puelche S.A.}
\item Distribuidora Puelche S.A. (2026d). \emph{Aclaraciones de la Licitación N.º TFEP-01/2026}.
\item LafroX. (2026). Subdocumentos 1 a 9 y 13, con los anexos y formularios citados, incluidos los Formularios T-13 y T-17.
\item International Electrotechnical Commission. (2018). \emph{IEC 60812:2018 Failure modes and effects analysis (FMEA and FMECA)}. IEC.
\item International Organization for Standardization. (2018). \emph{ISO 31000:2018 Risk management \textemdash{} Guidelines}. ISO.
\item Project Management Institute. (2017). \emph{La guía de los fundamentos para la dirección de proyectos (Guía del PMBOK®)} (6.ª ed.), capítulo 11. Project Management Institute.
\end{itemize}

\section*{Declaración de uso de IA}
\addcontentsline{toc}{section}{Declaración de uso de IA}

La tabla declara apoyo de IA conforme a las Aclaraciones §7.2. La revisión humana identifica quién efectivamente verificó cada parte y se consolida en el Formulario A-6.

\begingroup
\renewcommand{\thetable}{8.IA.1}
\begin{tablalafrox}{Declaración de uso de IA}{tab:sd8-cuerpo-7}{>{\hsize=0.98597\hsize\linewidth=\hsize\RaggedRight\arraybackslash}X >{\hsize=1.04345\hsize\linewidth=\hsize\RaggedRight\arraybackslash}X >{\hsize=1.51284\hsize\linewidth=\hsize\RaggedRight\arraybackslash}X >{\hsize=0.56762\hsize\linewidth=\hsize\RaggedRight\arraybackslash}X >{\hsize=0.71876\hsize\linewidth=\hsize\RaggedRight\arraybackslash}X >{\hsize=1.17136\hsize\linewidth=\hsize\RaggedRight\arraybackslash}X}{\cab{Sección} & \cab{Herramienta} & \cab{Finalidad del uso} & \cab{Nivel en texto} & \cab{Nivel en diagramas} & \cab{Revisión humana (quién y qué verificó)}}
Introducción & Codex & Conexión con operación, alcance y formularios & Alto & Ninguno & [[REVISIÓN HUMANA]] \\
8.1 Plan de riesgos & Codex; Claude Code & Roles, escalas, ciclo y cierre & Alto & Ninguno & [[REVISIÓN HUMANA]] \\
8.2 Identificación y Análisis de Riesgos & Codex; Claude Code & RBS, exposición cualitativa y valor esperado en HH & Alto & Medio & [[REVISIÓN HUMANA]] \\
8.3 Plan de Acción a Riesgos & Codex; Claude Code & Respuestas, contingencia y gestión & Alto & Ninguno & [[REVISIÓN HUMANA]] \\
Anexo 8.A & Codex; Claude Code & Ver declaración propia de ese anexo & Alto & Ninguno & [[REVISIÓN HUMANA]] \\
Anexo 8.B & Codex; Claude Code & Ver declaración propia de ese anexo & Alto & Ninguno & [[REVISIÓN HUMANA]] \\
Anexo 8.C & Codex; Claude Code & Ver declaración propia de ese anexo & Alto & Ninguno & [[REVISIÓN HUMANA]] \\
Anexo 8.D & Codex; Claude Code & Ver declaración propia de ese anexo & Alto & Ninguno & [[REVISIÓN HUMANA]] \\
Anexo 8.E & Codex; Claude Code & Ver declaración propia de ese anexo & Alto & Ninguno & [[REVISIÓN HUMANA]] \\
Anexo 8.F & Codex; Claude Code & Ver declaración propia de ese anexo & Alto & Ninguno & [[REVISIÓN HUMANA]] \\
Anexo 8.G & Codex & Ver declaración propia de ese anexo & Alto & Ninguno & [[REVISIÓN HUMANA]] \\
Formulario T-16 & Codex; Claude Code & Ver declaración propia del formulario & Alto & Ninguno & [[REVISIÓN HUMANA]] \\
\end{tablalafrox}
\endgroup